
\documentclass[11pt]{article}
\usepackage[margin=0.8in]{geometry}
\usepackage{amsmath,amssymb,booktabs,array,longtable}
\usepackage{hyperref}
\usepackage{enumitem}
\usepackage{graphicx}
\usepackage{float}
\setlist{nosep}

\begin{document}

\begin{center}
{\Large\textbf{Shiv Nadar University Chennai}}\\
\textbf{CS3807 -- Deep Learning Laboratory}\\
\textbf{Experiment 2}\\
{\large\textbf{Implementation of a Multi-Layer Perceptron (MLP) for Multi-Class Image Classification}}
\end{center}

\renewcommand{\arraystretch}{1.25}

\begin{tabular}{|p{3.8cm}|p{8cm}|p{2cm}|}
\hline
Degree \& Branch & B.Tech Artificial Intelligence \& Data Science & Semester V\\ \hline
Subject Code \& Name & CS3807 -- Deep Learning Laboratory & AY: 2026--27\\ \hline
\end{tabular}

\section*{1. Objective}
Implement an MLP using TensorFlow/Keras on the Fashion-MNIST dataset. Learn image preprocessing, flattening, model construction, training, evaluation, and automated hyperparameter optimization.

\section*{2. Background Theory}
An MLP consists of an input layer, hidden layers and an output layer.

\[
z^{(l)}=W^{(l)}a^{(l-1)}+b^{(l)},\qquad
a^{(l)}=f(z^{(l)})
\]

Output layer:
\[
\hat y=\mathrm{Softmax}(z)
\]

Loss:
\[
L=-\sum_i y_i\log(\hat y_i)
\]

Recommended architecture:

\begin{center}
784 $\rightarrow$ Dense(128, ReLU) $\rightarrow$ Dense(64, ReLU) $\rightarrow$ Dense(10, Softmax)
\end{center}

\section*{3. Dataset}
Dataset: Fashion-MNIST

Training Images: 60,000 \\
Testing Images: 10,000 \\
Classes: 10 \\
Image Size: $28\times28$

\section*{4. Image Preprocessing}

Fashion-MNIST images are stored as $28\times28$ matrices. Dense layers require a one-dimensional feature vector.

\[
28\times28=784
\]

Example:

\[
\begin{bmatrix}
10&20\\30&40
\end{bmatrix}
\rightarrow
[10,20,30,40]
\]

Perform the following preprocessing:

\begin{enumerate}
\item Load Fashion-MNIST.
\item Display sample images.
\item Flatten images.
\item Normalize pixels to [0,1].
\item Convert labels into one-hot vectors.
\end{enumerate}

\section*{5. Source Code}

The complete implementation (dataset loading, preprocessing, baseline MLP, training, evaluation, and hyperparameter optimization) is available at:

\begin{center}
\url{https://github.com/Jai-Abhishek/DL-Lab}
\end{center}

\section*{6. Hyperparameter Optimization}

Instead of manually selecting hyperparameters, students shall perform automated hyperparameter optimization using either \textbf{GridSearchCV} or \textbf{RandomizedSearchCV} (recommended) together with the SciKeras wrapper.

Optimize the following search space.

\begin{center}
\begin{tabular}{|p{5cm}|p{7cm}|}
\hline
Hyperparameter & Candidate Values\\
\hline
Hidden Layers & 1, 2, 3\\
Hidden Neurons & 32, 64, 128, 256\\
Learning Rate & 0.1, 0.01, 0.001\\
Batch Size & 16, 32, 64, 128\\
Epochs & 10, 20, 30\\
Optimizer & SGD, Adam, RMSProp\\
Activation Function & ReLU, Tanh, Sigmoid\\
Dropout Rate (Optional) & 0.0, 0.2, 0.5\\
\hline
\end{tabular}
\end{center}

\subsection*{Tasks}

\begin{enumerate}
\item Build a baseline MLP model.
\item Define the hyperparameter search space.
\item Perform Grid Search or Randomized Search using 5-fold cross-validation.
\item Record the best hyperparameter combination.
\item Retrain the model using the optimized hyperparameters.
\item Evaluate the optimized model on the testing dataset.
\item Compare the optimized model with the baseline model.
\end{enumerate}

\subsection*{Results of Hyperparameter Optimization}

A \texttt{RandomizedSearchCV} (5-fold cross-validation, 10 sampled iterations) was performed using the SciKeras wrapper on a 10,000-sample subset of the training data, over the search space defined above. The search returned the following best configuration:

\begin{center}
\begin{tabular}{|p{6cm}|p{6cm}|}
\hline
Hidden Layers & 3\\ \hline
Hidden Neurons & 256\\ \hline
Learning Rate & 0.001\\ \hline
Batch Size & 32\\ \hline
Optimizer & Adam\\ \hline
Activation Function & Tanh\\ \hline
Epochs & 20\\ \hline
Dropout Rate & 0.2\\ \hline
Best Cross-validation Accuracy & 0.8504\\ \hline
\end{tabular}
\end{center}

This optimized configuration was then retrained on the full training set and evaluated on the held-out test set, achieving a testing accuracy of 0.8783, compared to 0.8758 for the baseline model (see Section 8 for the full comparison).

\section*{7. Mandatory Plots}

The plots below were generated from the notebook. A short inference is provided under each plot.

\begin{figure}[H]
\centering
\includegraphics[width=0.85\textwidth]{figures/sample_images.png}
\caption{Sample Images from Fashion-MNIST}
\end{figure}
The ten sample images confirm that the class labels correctly correspond to their respective garment categories (e.g., trouser, coat, sneaker). The low-resolution ($28\times28$) grayscale format shows that fine details such as texture are barely distinguishable, which explains why visually similar classes are harder to separate.

\begin{figure}[H]
\centering
\includegraphics[width=0.7\textwidth]{figures/class_distribution.png}
\caption{Class Distribution in Training Set}
\end{figure}
All ten classes contain exactly 6,000 training samples each, so the training set is perfectly balanced. This removes class imbalance as a possible cause of the lower per-class performance observed later for classes such as Shirt.

\begin{figure}[H]
\centering
\includegraphics[width=0.6\textwidth]{figures/train_accuracy.png}
\caption{Training Accuracy vs Epoch (Baseline)}
\end{figure}
Training accuracy increases steadily from about 0.82 to 0.94 over 20 epochs, with the steepest gains occurring in the first five epochs. The smooth, monotonic curve indicates that the model is learning consistently without instability in the optimization process.

\begin{figure}[H]
\centering
\includegraphics[width=0.6\textwidth]{figures/val_accuracy.png}
\caption{Validation Accuracy vs Epoch (Baseline)}
\end{figure}
Validation accuracy rises quickly in the first few epochs and then plateaus around 0.88, fluctuating slightly for the remaining epochs. This plateauing, while training accuracy keeps rising, is an early indicator of mild overfitting.

\begin{figure}[H]
\centering
\includegraphics[width=0.6\textwidth]{figures/train_loss.png}
\caption{Training Loss vs Epoch (Baseline)}
\end{figure}
Training loss decreases smoothly and continuously from roughly 0.51 to 0.16 across the 20 epochs. The consistently decreasing curve shows the model is fitting the training data increasingly well without any signs of divergence.

\begin{figure}[H]
\centering
\includegraphics[width=0.6\textwidth]{figures/val_loss.png}
\caption{Validation Loss vs Epoch (Baseline)}
\end{figure}
Validation loss decreases initially but starts rising again after around epoch 8--10, while training loss keeps falling. This divergence between training and validation loss is a clear sign of overfitting in the later epochs of baseline training.

\begin{figure}[H]
\centering
\includegraphics[width=0.7\textwidth]{figures/confusion_matrix.png}
\caption{Confusion Matrix (Baseline Model)}
\end{figure}
Most classes are classified with high accuracy, particularly Trouser, Sandal, and Ankle boot, which lie almost entirely on the diagonal. The largest confusions occur between Shirt and T-shirt/top, Pullover, and Coat, since these upper-body garments share similar silhouettes at $28\times28$ resolution.

\begin{figure}[H]
\centering
\includegraphics[width=0.8\textwidth]{figures/hp_search.png}
\caption{Hyperparameter Search Iterations vs Mean CV Score}
\end{figure}
The mean cross-validation score varies considerably across the 10 randomly sampled configurations, ranging from roughly 0.65 to 0.85. This wide spread confirms that hyperparameter choice has a substantial effect on model performance and justifies the use of a systematic search rather than manual tuning.

\begin{figure}[H]
\centering
\includegraphics[width=0.55\textwidth]{figures/model_comparison.png}
\caption{Baseline vs Optimized Model Accuracy Comparison}
\end{figure}
The optimized model achieves a marginally higher test accuracy (0.8783) than the baseline model (0.8758). The improvement is small, suggesting the manually chosen baseline architecture was already a reasonably good configuration for this dataset.

\section*{8. Results}

\textbf{Best Hyperparameters}

\begin{tabular}{|p{6cm}|p{6cm}|}
\hline
Hidden Layers & 3\\\hline
Hidden Neurons & 256\\\hline
Learning Rate & 0.001\\\hline
Batch Size & 32\\\hline
Optimizer & Adam\\\hline
Activation Function & Tanh\\\hline
Epochs & 20\\\hline
Dropout & 0.2\\\hline
Cross-validation Accuracy & 0.8504\\\hline
Testing Accuracy & 0.8783\\\hline
\end{tabular}

\vspace{3mm}

\textbf{Performance Comparison}

\begin{tabular}{|p{4cm}|c|c|}
\hline
Metric & Baseline & Optimized\\
\hline
Accuracy&0.8758&0.8783\\
Precision&0.8772&0.8800\\
Recall&0.8758&0.8783\\
F1-score&0.8759&0.8785\\
Training Time (s)&132.81&123.85\\
\hline
\end{tabular}

\section*{9. Discussion and Conclusion}

\textbf{Which optimization method was used?}
A Randomized Search with 5-fold cross-validation (\texttt{RandomizedSearchCV}, 10 sampled configurations) was used, implemented via the SciKeras wrapper around the Keras \texttt{Sequential} model.

\textbf{Which hyperparameters were selected?}
The best configuration found was 3 hidden layers of 256 neurons each, Tanh activation, a dropout rate of 0.2, the Adam optimizer with a learning rate of 0.001, batch size 32, and 20 training epochs.

\textbf{Did optimization improve performance?}
Yes, but only marginally. Test accuracy improved from 0.8758 (baseline) to 0.8783 (optimized), a gain of about 0.25 percentage points, with similar small improvements in precision, recall and F1-score. The optimized model also trained slightly faster (123.85s vs 132.81s).

\textbf{Which hyperparameter had the greatest impact?}
Based on the spread of scores across search iterations (Plot 8), the number of hidden layers/neurons and the choice of optimizer and learning rate had the largest effect on cross-validation accuracy; configurations with a high learning rate (e.g., 0.1) or SGD optimizer produced noticeably lower scores than Adam/RMSProp with a small learning rate.

\textbf{Compare Grid Search and Randomized Search.}
Grid Search exhaustively evaluates every combination in the search space, which is guaranteed to find the best combination within the grid but is computationally expensive as the space grows. Randomized Search instead samples a fixed number of combinations, making it far cheaper to run on a large search space (as used here) while still typically finding a near-optimal configuration.

\textbf{Recommend the best MLP configuration.}
Given the small accuracy gain relative to the added complexity, the optimized architecture (3 hidden layers of 256 neurons, Tanh activation, dropout 0.2, Adam optimizer, learning rate 0.001) is recommended when the extra training cost is acceptable. For quick baselines or resource-constrained settings, the simpler 2-layer (128, 64) ReLU architecture remains a reasonable choice, since it achieves nearly the same accuracy at lower computational cost.

In conclusion, both the baseline and the hyperparameter-optimized MLP classify Fashion-MNIST images with approximately 88\% test accuracy. Automated hyperparameter search using \texttt{RandomizedSearchCV} was able to identify a deeper, regularized architecture that offered a small but consistent improvement over the manually chosen baseline, confirming the value of systematic hyperparameter optimization even when the achievable gains are modest.

\section*{10. References}
\begin{enumerate}
\item Goodfellow et al., \emph{Deep Learning}.
\item Bishop, \emph{Pattern Recognition and Machine Learning}.
\item Haykin, \emph{Neural Networks and Learning Machines}.
\item Fashion-MNIST Dataset.
\item TensorFlow/Keras Documentation.
\end{enumerate}

\end{document}
